% PN board factory signed-firmware flashing workflow -- serpentine TikZ flowchart.
% Engine: XeLaTeX (xeCJK).  Build with:  make
\documentclass[a4paper]{article}
\usepackage[a4paper,margin=1cm]{geometry}
\usepackage{xeCJK}
\usepackage{tikz}
\usetikzlibrary{arrows.meta, positioning, shapes.geometric}

% CJK fonts -- change here if the build host has a different font set installed.
\setCJKmainfont{Noto Serif CJK SC}
\setCJKsansfont{Noto Sans CJK SC}

\pagestyle{empty}

\begin{document}

\begin{center}
{\Large \textbf{PN 板卡出厂烧录签名固件流程}}
\end{center}
\vspace{4mm}

\begin{tikzpicture}[
  font=\footnotesize,
  >=Stealth,
  % node distance = <vertical gap> and <horizontal gap> between node borders
  node distance=21mm and 7mm,
  every node/.style={align=center},
  box/.style={
    draw, rounded corners=2pt, thick,
    minimum width=3.9cm, minimum height=1.5cm,
    text width=3.6cm, align=center, inner sep=2pt
  },
  mode/.style={box, fill=blue!12, draw=blue!60!black},
  power/.style={box, fill=green!12, draw=green!50!black},
  poweroff/.style={box, fill=gray!18, draw=gray!60!black},
  op/.style={box, fill=orange!15, draw=orange!70!black},
  % Dashed frame = mass-production / pentest stages only (colours unchanged).
  prodonly/.style={dashed},
  arr/.style={-{Stealth[length=2.5mm]}, very thick}
]

% ---------- 第一行：1 -> 2 -> 3 -> 4 (从左到右) ----------
\node[mode]                    (n1) {\textbf{1.} 工装切换启动模式为 SCI 模式};
\node[power,    right=of n1]   (n2) {\textbf{2.} 板卡上电};
\node[op,       right=of n2]   (n3) {\textbf{3.} 通过串口烧录 OTP RPKH\footnotemark};
\node[poweroff, right=of n3]   (n4) {\textbf{4.} 板卡断电};

% ---------- 第二行：5 -> 6 -> 7 -> 8 (从右到左) ----------
\node[mode,     below=of n4]   (n5) {\textbf{5.} 工装切换启动模式为 xSPI0 模式};
\node[power,    left=of n5]    (n6) {\textbf{6.} 板卡上电};
\node[op,       left=of n6]    (n7) {\textbf{7.} 通过 JTAG 烧录签名固件};
\node[poweroff, left=of n7]    (n8) {\textbf{8.} 板卡断电};

% ---------- 第三行：9 -> 10 -> 11 -> 12 (从左到右) ----------
\node[power,    below=of n8]   (n9)  {\textbf{9.} 板卡上电\\(保持 xSPI0 启动模式)};
\node[op,       right=of n9]   (n10) {\textbf{10.} 通过固件烧录加密芯片};
\node[poweroff, right=of n10]  (n11) {\textbf{11.} 板卡断电};
\node[mode,     prodonly, right=of n11]  (n12) {\textbf{12.} 工装切换启动模式为 SCI 模式};

% ---------- 第四行：13 -> 14 -> 15 -> 16 (从右到左) ----------
\node[power,    prodonly, below=of n12]  (n13) {\textbf{13.} 板卡上电};
\node[op,       prodonly, left=of n13]   (n14) {\textbf{14.} 永久禁止 JTAG};
\node[op,       prodonly, left=of n14]   (n15) {\textbf{15.} 检查 JTAG 状态};
\node[poweroff, prodonly, left=of n15]   (n16) {\textbf{16.} 板卡断电};

% ---------- 第五行：17 -> 18 -> 19 (从左到右) ----------
\node[mode,     prodonly, below=of n16]  (n17) {\textbf{17.} 工装切换启动模式为 SCI 模式};
\node[power,    prodonly, right=of n17]  (n18) {\textbf{18.} 板卡上电};
\node[op,       prodonly, right=of n18]  (n19) {\textbf{19.} 通过串口永久封锁 SCI 启动模式};

% ---------- 箭头连接 ----------
% 第一行 (向右)
\draw[arr] (n1) -- (n2);
\draw[arr] (n2) -- (n3);
\draw[arr] (n3) -- (n4);
% 行间转折 (右侧向下)
\draw[arr] (n4) -- (n5);
% 第二行 (向左)
\draw[arr] (n5) -- (n6);
\draw[arr] (n6) -- (n7);
\draw[arr] (n7) -- (n8);
% 行间转折 (左侧向下)
\draw[arr] (n8) -- (n9);
% 第三行 (向右)
\draw[arr] (n9)  -- (n10);
\draw[arr] (n10) -- (n11);
\draw[arr] (n11) -- (n12);
% 行间转折 (右侧向下)
\draw[arr] (n12) -- (n13);
% 第四行 (向左)
\draw[arr] (n13) -- (n14);
\draw[arr] (n14) -- (n15);
\draw[arr] (n15) -- (n16);
% 行间转折 (左侧向下)
\draw[arr] (n16) -- (n17);
% 第五行 (向右)
\draw[arr] (n17) -- (n18);
\draw[arr] (n18) -- (n19);

\end{tikzpicture}

% ---------- Note explaining the dashed frames ----------
\vspace{5mm}
\begin{center}
\footnotesize
\tikz[baseline=-0.55ex]{\node[draw, dashed, thick, rounded corners=2pt,
  minimum width=0.5cm, minimum height=0.4cm, inner sep=0pt] {};}\;
\textbf{注：}虚线边框的步骤（第 12 步及以后）仅量产阶段和 Pentest 送测阶段适用，试产阶段不执行。
\end{center}

\vspace{4mm}
\begin{center}
\begin{tikzpicture}[font=\footnotesize, every node/.style={align=left}]
  \node[draw, rounded corners=2pt, fill=blue!12, draw=blue!60!black, minimum width=0.5cm, minimum height=0.4cm] (l1) at (0,0) {};
  \node[right=1mm of l1] {启动模式切换};
  \node[draw, rounded corners=2pt, fill=green!12, draw=green!50!black, minimum width=0.5cm, minimum height=0.4cm] (l2) at (4.2,0) {};
  \node[right=1mm of l2] {板卡上电};
  \node[draw, rounded corners=2pt, fill=gray!18, draw=gray!60!black, minimum width=0.5cm, minimum height=0.4cm] (l3) at (7.6,0) {};
  \node[right=1mm of l3] {板卡断电};
  \node[draw, rounded corners=2pt, fill=orange!15, draw=orange!70!black, minimum width=0.5cm, minimum height=0.4cm] (l4) at (11.0,0) {};
  \node[right=1mm of l4] {烧录 / 操作步骤};
\end{tikzpicture}
\end{center}

\footnotetext{RPKH：Root Public Key Hash，芯片安全启动信任根公钥哈希，烧录后不可更改，请务必核对密钥哈希值后再执行。}

\end{document}
